% nam-mf.tex -- TUGboat article (plain TeX, tugboat.sty).
% Typeset with:  luatex nam-mf.tex
% Lua\TeX\ is needed only for the few Korean words; the Korean font
% is UnBatang from TeX Live (unfonts-core), loaded by file name.
\input tugboat.sty
\input luaotfload.sty
\StretchyTenPointSpacing
\StretchyNinePointSpacing

\font\kofont={file:UnBatang.ttf} at 9.5pt
\def\ko#1{{\kofont #1}}
\def\GWEB{{\tt GWEB}\spacefactor1000 }
\def\CWEB{{\tt CWEB}\spacefactor1000 }
\def\LuaTeX{Lua\TeX}
\def\MFTP{{\sl \slMF: The Program\/}}
\def\id#1{\hbox{\it #1\/}}
\def\sect#1{$\langle\,$#1$\,\rangle$}
\def\url#1{{\tt #1}}
\def\image#1#2{\saveimageresource width #1 {#2}%
  \useimageresource\lastsavedimageresourceindex}
\def\figcaption#1#2{\smallskip
  {\ninepoint\noindent{\bf #1}\enspace #2\par}}

\title * \MFTP\ rewritten with Korean \GWEB\ *
\author * Soojin Nam *
\address * South Korea *
\netaddress * jsunam@gmail.com\\ https://sjnam.github.io *

\article

\subhead * Abstract *
I rewrote Knuth's \MFTP\ ({\tt mf.web}, version 2.71828182) in full
as a literate program. The prose is in Korean, the code is in Go, and
the tool is \GWEB, which I wrote. The result passes the \acro{TRAP}
test and produces eleven Computer Modern fonts identical, byte for byte,
to those made by the \MF\ of \TeX\ Live. I did this work together with
Claude, an \acro{AI} model from Anthropic. This article tells how
\GWEB\ came about and how this work improved it, the problems of
translating into a human language and into a programming language, and
a small error in the original that I found along the way.

\head * 1\quad Introduction *

I am a retired programmer who has used \TeX\ for more than thirty
years. Since retiring I have had plenty of time, and that time, along
with being a fan of Knuth, made this work possible. I had two reasons
for starting it.

First, I wanted to learn from a large program written by a master: how
to build a program, and how to turn algorithms into concrete code.
\MF\ seemed to have even more interesting algorithms than \TeX. To
read such a masterpiece again and again and understand it more deeply,
translating it into Korean,
my native language, was the natural thing for me to do. One cannot
translate a sentence without understanding it, so translation is also
the most thorough kind of reading.

Second, I wanted to see whether my \GWEB\ works well on a program as
large and complex as {\tt mf.web}, and to improve \GWEB\ in the
process. To test a tool, a large, well-structured program is better
than a toy program, so choosing Knuth's programs was natural. There is
no better test than real use.

Before this, I translated \TP\ in the same way
(tex-the-program~[9]); this work (metafont-the-program~[8]) applies
that experience to \MF. As far as I know, {\tt tex.web} has been
translated into many programming languages, but translations of
{\tt mf.web} are rare, and neither has been translated into another
human language. In that sense this translation is unusual in two ways.

The result, {\tt mf.w}, has 29,153 lines, and the Go code that
{\tt gtangle} extracts from it has 23,830 lines. It follows the 52
parts of the original and is 856 pages when typeset. The original has
1,215 sections and the translation 1,353 (section~4 explains why). In
this article a \S\ number is that of the original unless stated
otherwise; when I need a section number of the translation, I write it
as ``translation \S161''.

\head * 2\quad \GWEB\ *

\subhead * Beginnings *
Some time ago I ported Knuth's {\tt advent.w}~[3] and
{\sl The Stanford GraphBase\/}~[4] to Go. I ported them as ordinary Go
programs, not as literate programs. When I sent Knuth an error I had
found in his work, at his address for bug reports, I also told him
about these ports. He regretted that they were not literate programs
and suggested that I write a tool for literate programming in Go. The
name \GWEB\ first appears in that reply, which I received on June~18,
2026. The first commit of \GWEB\ was made the next day.

{\narrower\smallskip\noindent
About porting SGB and ADVENTURE to Go: I'm very glad that you have
enjoyed these algorithms. But I'm disappointed that you seem to have
eliminated the ``literate'' aspect! Is there not a tool for literate
programming in Go? If not, why not write GWEB? \dots\par
\rightline{Best wishes, Don}\smallskip}

The \GWEB\ that came of it~[6] is now introduced on Knuth's literate
programming page~[5] in one line: ``Soojin Nam's \GWEB, an analog of
\CWEB\ for Go instead of C''.

\subhead * Design *
As the name says, \GWEB\ follows \CWEB. {\tt gtangle} extracts the Go
code and {\tt gweave} makes a \TeX\ document. Change files, {\tt @i}
includes, the index and the list of section names, even the page
layout follow \CWEB. Three things differ because of Go.

\list\raggedright
\item {\tt gtangle} formats the assembled program with {\tt gofmt},
and it always adds {\tt //line} directives. Like the {\tt \#line} of
\CWEB, they make the compiler's error messages and panic traces point
to lines of the {\tt .w} file instead of the {\tt .go} file.
\item {\tt gweave} does not rebuild the layout of the code from a
grammar as \CWEB\ does; it follows the spacing of the source. Go
formatted by {\tt gofmt} already carries its grammar in its spacing
(|*T| versus |a * b|, |a*b + c|, and so on).
\item A Korean document needs only one line, \cs{input kotexgweb}, in
its limbo. It is typeset with \LuaTeX\ and ko.\TeX\ (luatexko), and
the words that {\tt gweave} generates, such as ``Names of the
sections'', become Korean too.
\endlist

\GWEB\ is written in \GWEB\ itself, and rebuilding it gives a
byte-identical result. Besides Knuth's two programs, I have written
about thirty literate Go programs in Korean with it~[7].

\subhead * What {\tt mf.web} taught \GWEB\ *
Knuth's \WEB\ sometimes puts a section name in the middle of an
expression, as in {\tt if @<Character \char`\|k\char`\|\ cannot be
printed@> then}. In
C this causes no trouble, but in \GWEB\ such code did not compile.
{\tt gtangle} put newlines before and after an expanded section, and
Go automatically inserts a semicolon at a newline after an identifier,
a literal, or a closing bracket, ending the statement. The opening
brace after the expanded condition became a syntax error. So in both
translations I avoided the problem by storing the condition in a
boolean variable first and using that in the |if|: four places in the
\TeX\ translation, two in the \MF\ translation.

This time I removed the workaround. Now {\tt gtangle} expands a section
reference that has other code on the same line in place, without
newlines. The position information is attached in the form
|/*line file:line*/|, which Go accepts even in the middle of a line.
Then error messages point, inside the expanded text, to the line where
the section is defined, and after it to the original line:

\verbatim
if /*line mf.w:1406*/ k < ' ' ||
   k > '~' /*line mf.w:1369*/ {
\endverbatim

\noindent (In {\tt gtangle}'s output this is one line.) I released this
change as \GWEB\ 0.10.4 and then restored all six places to the
original form. The tests of both translations still passed. In the same
release the type style of constants was changed to follow the rule of
{\tt cweave}: a name with no lowercase letters is set in typewriter
type. I had expected that translating a large program would expose gaps
in the tool, and that is exactly what happened.

\head * 3\quad Translating into Korean *

\subhead * What was translated and what was not *
I translated the prose, the section names, and the comments into
Korean. The identifiers, the index entries, and the error and help
messages that \MF\ prints remain in English. If the messages were
translated, the output could not be compared byte by byte with the
\acro{TRAP} test and with \TeX\ Live. I kept the index in English so
that it is easy to use side by side with the book.

\subhead * Section names *
Knuth's section names are usually English imperatives or short
statements. In Korean the verb comes at the end of the sentence, so I
translated section names into statements that end with the verb. For
example, \sect{Character $k$ cannot be printed} becomes
\sect{\ko{문자} $k$\ko{를 찍을 수 없다}}. Within the code, a section
name reads as a short Korean sentence of one line.

\subhead * Identifiers and particles *
A Korean particle changes its form depending on the sound that ends
the preceding word: $k$\ko{를}, \id{mMagic}\ko{과},
\id{futurePen}\ko{은}. (The letter $k$ is read ``kay'', ending in a
vowel, so it takes \ko{를}; \id{futurePen} ends in a consonant, so it
takes \ko{은}.) So to attach a particle to an English identifier, one
must decide how the identifier is pronounced. Also, a Korean sentence
that begins with an identifier or a formula reads very awkwardly. So
instead of ``\id{fraction}\ko{은} \dots'' I wrote ``\ko{형}
\id{fraction}\ko{은} \dots'' (``the type \id{fraction} \dots''),
putting in front a Korean noun that names it. After finishing the
translation I searched again for sentences that broke this rule. There
were about 330, and I fixed them all.

\subhead * Terms *
Most \MF\ terms have no established Korean translation. Where I could,
I chose plain native Korean words (table~1).

\medskip
\centerline{\ninepoint{\bf Table 1.}\enspace Some translated terms.}
\smallskip
\centerline{\ninepoint\vbox{\halign{#\hfil\tabskip=1.5em
  &\ko{#}\hfil\tabskip=0pt\cr
scaled (integer)&눈금 정수\cr
word&낱말\cr
halfword, quarterword&하프워드, 쿼터워드\cr
base file, transcript file&기반 파일, 기록 파일\cr
octant&팔분면\cr
edge structure&가장자리 구조\cr
mock curvature&모조 곡률\cr
dependency list&종속 목록\cr
granularity&낱알 크기\cr
pseudoprinting&흉내 찍기\cr
Stern--Brocot wreath&스턴–브로코 화환\cr}}}
\medskip

\subhead * Knuth's voice *
Knuth's writing has much humor. I translated ``That which can be
created can be destroyed'' (\S487) as ``\ko{만들어질 수 있는 것은
부서질 수 있다}'', and ``The former chastity has vanished!''
(\S1203) as ``\ko{이전의 순결은 사라졌다!}''. If I translate only the
meaning, the joke is lost; if I keep the flavor of the words, the
meaning blurs. There were many places where I struggled between the two.

\head * 4\quad Translating into Go *

The goal was to remove the smell of Pascal and write in the Go way,
without changing Knuth's arithmetic or data structures, because the
output had to be the same down to the byte. Table~2 shows the main
correspondences.

\medskip
\centerline{\ninepoint{\bf Table 2.}\enspace Pascal in the original,
Go in the translation.}
\smallskip
{\ninepoint\raggedright
\dimen0=.47\hsize
\halign to\hsize{\vtop{\hsize=\dimen0 \noindent#\strut}\tabskip=0pt plus1fil
  &\vtop{\hsize=\dimen0 \noindent#\strut}\tabskip=0pt\cr
|goto end_of_MF| from another procedure&
  |panic| with the label as value; |recover| in the main program\cr
\noalign{\smallskip}
variant record |memory_word|&a struct with separate fields
  (36~bytes)\cr
\noalign{\smallskip}
|@d half(#)==(#) div 2|&|func half(x int32) int32|
  |{ return x / 2 }|\cr
\noalign{\smallskip}
string constants preloaded in the string pool&the list
  |poolLiterals|; a checksum in the base file\cr
\noalign{\smallskip}
|forward| declarations&not needed; all functions are collected in one
  section, \sect{\ko{함수들}} (Functions)\cr
\noalign{\smallskip}
local labels |done|, |found|&mostly |break|, |continue|, labeled
  loops\cr
\noalign{\smallskip}
|stat| \dots\ |tats|, |debug| \dots\ |gubed|&|if stat {...}|,
  |if debug {...}| (constant conditions)\cr}}
\medskip

\noindent A few need explanation.

\subhead * A {\tt goto} out of a function *
When a fatal error occurs, \MF\ jumps from |jump_out| straight to a
label in the main program. Go's |goto| cannot leave a function, so the
translation raises a |panic| with the label as its value, and a
function that the main program has deferred catches it with |recover|.

\subhead * The memory word *
The original |memory_word| is a variant record that overlays one
integer, two halfwords, or four quarterwords in one word. Go has no
such record. \MF\ never reads a value through a different variant from
the one it was written with, so instead of using |unsafe| I chose a
struct with separate fields. The price is size. One word became
36~bytes, so the base file is about four times as large as that of
\TeX\ Live (299\thinspace KB versus 78\thinspace KB), and computation
became slower.

\subhead * Integer division *
Knuth recommends that |half| be changed to a shift for speed. But Go's
|/| truncates toward zero, while |x>>1| rounds toward $-\infty$
(|-3/2| is $-1$; |-3>>1| is $-2$). The log of the \acro{TRAP} test was
made with the former arithmetic, so I kept the division. The Go
compiler turns |x/2| into a few shifts anyway, so nothing is lost.

\subhead * Functions and sections *
Procedures that the original declared in advance with |forward| are
defined in place, and all functions are appended to one section,
\sect{\ko{함수들}}. Most local labels of the original became |break|,
|continue|, or |if|--|else|, and at each such place the prose notes
what the original did. The number of sections grew because of
typesetting. {\tt gwebmac} breaks pages mainly at section boundaries,
so a section longer than a page causes an Underfull warning. So I split
long sections into named sections. The typesetting produces no
warnings.

\subhead * Performance *
Making cmr10 at 600\thinspace dpi took 0.18 seconds with the
translation and 0.13 seconds with {\tt mf-nowin} of \TeX\ Live (Apple
M1 Max). Pure computation is a little more than twice as slow in the
translation (probably because of the 36-byte word and array bounds
checks), but the translation starts in 0.01 seconds, while
{\tt mf-nowin} spends about 0.1 seconds initializing kpathsea. The
first version took 1.4 seconds, because it ran {\tt kpsewhich}
separately each time it looked for an input file. When I made it scan
the input directories only once and build a table, the time became
0.18 seconds.

\head * 5\quad Verification *

\subhead * The \acro{TRAP} test *
The {\tt trap.72270gf} made by the translation is identical to the
original's byte for byte, and {\tt trap.pl} and {\tt trap.typ} have no
differing lines. The only differences in the log files are the banner
and the string pool statistics, both of which the \acro{TRAP} manual
allows.

\subhead * Comparison with \TeX\ Live *
I made eleven Computer Modern fonts (cmr10, cmbx12, cmti10, cmmi10,
cmsy10, cmex10, cmtt10, cmss10, cmcsc10, cmfib8, logo10) at
200\thinspace dpi and 600\thinspace dpi and compared them with the
output of {\tt mf-nowin}. In all 22 cases the \TFM, \acro{GF}, and log
files were the same. The original |fix_date_and_time| sets the date to
noon on July~4, 1776, if it cannot get the system date; the translation
deliberately always uses that date. Then the outputs can be compared
directly, except for the date comment in the \acro{GF} preamble.

\head * 6\quad What careful reading found *

\subhead * The {\tt link(p)} of \S587 *
The paragraph that describes dependency lists reads:

{\narrower\smallskip\noindent
If $q=\id{dep\_list}(p)$ points to this list, and if $k>0$, then
$\id{value}(q)=\alpha_1$ (which is a \id{fraction}); $\id{info}(q)$
points to the location of $v_1$; and $\id{link}(p)$ points to the
dependency list $\alpha_2v_2+\cdots+\alpha_kv_k+\beta$.
\smallskip\par}

The final $\id{link}(p)$ should be $\id{link}(q)$. The sentence lists
three fields of the first node, which it has just named~$q$, and the
first two are $\id{value}(q)$ and $\id{info}(q)$. The variable $p$ is
the dependent variable itself, and its \id{link} is used to connect the
variable structure (\S228ff.). The code that traverses the list
(\id{print\_dependency} in \S589 and others) also goes to the next term
through the \id{link} of a list node, and there the list node is called
$p$. Perhaps that name slipped into the paragraph. The error was still
in the 2021 version and was not in Knuth's errata. When I reported it
to tex-k, Karl Berry confirmed it and added it to the list to be sent
to Knuth.

\subhead * Theoretical 32-bit overflows *
The translation runs in a language that does not catch integer
overflow, so I examined, one by one, the places where overflow could
occur. Knuth checked almost every place in advance, but in three places
overflow can occur when a number of 32768 or more is to be represented
as a scaled integer (an integer times $2^{16}$).

\list[\numbered]\raggedright
\item {\bf The lower threshold of |cull| (\S1074).} The original rounds
the lower end up as |floor_unscaled(cur_x+unity-1)|. If a computed $a$
is 32767.5 or more, this addition overflows, the lower end becomes
$-32768$, pixels of weight~0 are kept too, and the picture is ruined.
In a version translated literally from the original, a picture culled
with $a=32767.5$ and $b=32767.9$ had a |totalweight| of $-0.1875$
instead of~0. The translation computes the ceiling as
$\lceil x\rceil=-\lfloor -x\rfloor$ instead. For values that do not
overflow the two formulas agree, so the result of the \acro{TRAP} test
does not change.
\item {\bf The window number of |display| (\S1071).} Converting an
out-of-range window number back with |cur_exp*unity| to show it in an
error message can overflow. Only the number in the message is wrong.
\item {\bf {\tt xoffset} and {\tt yoffset} (\S1165).} If an offset is
32767.5 or more, a wrong value goes into the |yyy| command of the
\acro{GF} file. The shapes of the characters are unaffected (the upper
part of figure~A).
\endlist

None of these shows up in real font work. The weight of a single edge
is at most $\pm3$, so for a pixel to reach weight 32768 one would have
to overlay more than ten thousand fills; an offset of 32768 pixels is
more than 1.4\thinspace m even at 600\thinspace dpi. Moreover, a
constant greater than 4096 is already cut off when it is scanned, so
such values arise only when computed on purpose. So I call these
``theoretical overflows'' rather than ``bugs'', and I did not report
them separately to tex-k. The people who look after the reports to
Knuth will also read this article, so I think writing them down here is
enough. The translation fixes the first and adds explanations for the
other two.

\head * 7\quad Figures *

Knuth put almost no figures in \MFTP. Instead he explains things very
precisely in words, and sometimes suggests that the reader draw a
picture with a pencil. In the translation I took up that suggestion and
drew 26 figures. They are drawn during typesetting with luamplib: the
\MP\ code sits in the {\tt .w} file itself, and since the
|btex| \dots\ |etex| labels are typeset by the document's own \TeX,
Korean labels automatically match the text font. Here are a few.

\list\raggedright
\item {\bf Octants and Gray code (\S139, translation \S161).} Writing
the original's codes in the eight octants shows that neighboring
octants differ in only one bit (figure~C).
\item {\bf The data structure of variables (\S229, translation
\S286).} Knuth explains the structure for |x|, |x.a|, |x[]b|, |x5|,
and |x20b| in words and says that the reader should draw a picture with
a pencil. The translation adds a section right after it and draws that
picture (figure~B).
\item {\bf Hobby's velocity function (\S116), de Casteljau subdivision
(\S303), Hobby's equivalence classes (\S525).} The approximation of
curves, the subdivision of Bernstein polynomials, and the number theory
behind elliptical pens, drawn following the examples of the original.
\item {\bf \acro{GF} encoding (\S1167, translation \S1291).} For a
character shaped like a |T|, five pixels wide, the figure shows row by
row which commands encode it (the lower part of figure~A).
\endlist

\figure
\centerline{\image{\hsize}{fig-b-variables.pdf}}
\figcaption{Figure B.}{The example of \S229 drawn (translation \S286):
the structure of variables after the user has mentioned |x|, |x.a|,
|x[]b|, |x5|, and |x20b| and declared |x[]b| to be {\bf boolean}.
Solid arrows are \id{link}, bold arrows \id{attr\_head}, and dashed
arrows \id{subscr\_head}; \id{parent} pointers are not drawn. The row
that starts at node $q_1$ of the collective subscript |x[]| has the
same shape as the row that starts at node $r_1$ of the actual subscript
|x20|. The structure under a collective subscript serves as the model
for the structure under an actual subscript.}
\endfigure

\figure
\centerline{\image{.75\hsize}{fig-c-octants-en.pdf}}
\figcaption{Figure C.}{Octant names and the Gray code (\S139, translation
\S161). The top of each cell gives the name of the octant; the bottom
gives the code defined in the original (\id{first\_octant}${}=1$,
\id{second\_octant}${}=5$, \dots). The three bits of the code minus~1
mean \id{negate\_x}, \id{negate\_y}, and \id{switch\_x\_and\_y}. Here
the names are in English, but in the figure of the translation they
are written in Korean, from \ko{첫째} (first) to \ko{여덟째} (eighth).}
\endfigure

\head * 8\quad Working with Claude *

This translation is not my work alone. Because my English is poor, I
did a large part of the translation and testing together with Claude,
an \acro{AI} model from Anthropic, used as a coding assistant (Claude
Code). Here is who did what.

\subhead * What I did *
I wrote \GWEB, on which this work is based, and I set the direction
and the standards. I wrote down in advance the rules for the structure
of the literate program (unnamed code only in the skeleton of the first
section; all functions collected in one section, \sect{\ko{함수들}}),
rules for Korean sentences, and the standard for typesetting quality
(zero warnings). Then I read the results and pointed out what to fix:
that the index section of the original was missing, that letters in
figures overlapped lines, and so on. It was also I who remembered the
workaround in \GWEB\ and proposed fixing it in \GWEB.

\subhead * What Claude did *
Claude read the original and wrote the Korean prose and the Go code,
drew the \MP\ figures, and wrote and ran test scripts. Claude was also
the first to notice the error in \S587. When I asked whether it was
sure this was a typo, it cited the definitions in the paragraph and the
code as evidence, and searched Knuth's errata and \acro{TUG}'s list to
confirm that it had not yet been reported.

\subhead * Working together *
Above all, the work took far less time than I had expected: less than
a week. Alone it would have taken me well over a year, and that only
if I had kept at it steadily and persistently. Moreover, the careful
and thorough testing done in that short time is something I could
never have imagined doing alone.

Working with Claude was one surprise after another. Sometimes it makes
mistakes that even a child would not make. It sometimes uses commands
or macros that do not exist in \GWEB\ as if they did, but such
mistakes can be forgiven as charming. I am most grateful when it
explains code I do not understand, in more detail than a book could
ever hold. It knows what I do not know better than I do, and I can only
wonder how that is possible. It writes example code, and sometimes
draws pictures to help me understand. When it brings up a reference I
would never have thought of for a piece of code, connecting things
that seemed unrelated, I cannot help being amazed.

\subhead * What made it trustworthy *
Some mistakes could not be taken lightly. At first it left out the
index section, and when measuring performance it almost reported times
measured while \TeX\ Live could not find the test file. What caught
such errors were three checks: Knuth's \acro{TRAP} test, the
byte-by-byte comparison with \TeX\ Live, and my own reading of the
typeset result. In a literate program the document and the code are
one, so a person checks understanding by reading the document, and
tests check the code. I never imagined that I would benefit from
Knuth's \acro{TRAP} in this way forty years later.

\head * 9\quad Conclusion *

Rereading a frozen program without skipping a single section taught me
why it has lasted more than forty years. The error I found this time
was a single letter in the prose. And \GWEB, which began with Knuth's
suggestion, came a step closer to \CWEB\ while translating Knuth's
program. The \MF\ translation was written as a literate program from
the start, so it is a late reply to that suggestion.

The translations and \GWEB\ are on GitHub~[6, 8, 9]. The copyright of
the original {\tt mf.web} belongs to Knuth. Following the conditions of
the original, this translation does not use the name {\tt mf.web} and
is not part of Knuth's \MF\ package.

\newpage\onecol
\dimen0=\vsize \advance\dimen0 by -8\baselineskip
\centerline{\saveimageresource height\dimen0 {fig-a-p767-crop.pdf}%
  \useimageresource\lastsavedimageresourceindex}
\figcaption{Figure A.}{Page 767 of the translation (translation
\S\S1290--1291, original \S\S1166--1167). The upper section is the
prose that explains the theoretical overflow of \id{xoffset}
(section~6) and the Go code that sends the offsets to the \acro{GF}
file; the lower section is the \acro{GF} encoding figure introduced in
section~7 and the Go code that outputs that character. In the code,
$\langle\dots\rangle$ is a named section, and the number after it is
the number of that section.}
\newpage\twocol

\head * Acknowledgments *

This project was possible thanks to the Korean \TeX\ Users Group
(\acro{KTUG}), which has worked for many years to make Korean usable in
\TeX. In particular, without Dohyun Kim~[10], who develops and
maintains luatexko and luamplib, on which Korean \GWEB\ depends
heavily, this project would not exist. I also thank Karl Berry, who
took an interest in translating {\tt mf.web} into Korean \GWEB\ and
encouraged me to write for \TUB.

Above all, all of this began with Knuth's work. Most of my programming
has been inspired by his many works. I do not dare to think what words
of thanks would be fitting for him. I simply want to convey my sincere
feelings with a quiet ``thank you''.

\head * References *

{\ninepoint\raggedright\frenchspacing
\def\bib#1 {\par\hangindent=1.8em\noindent\hbox to1.8em{[#1]\hfil}}
\def\s{\slash}
\bib1 Donald E. Knuth. \MFTP. Computers \& Typesetting, Volume~D.
  Addison-Wesley, 1986.
\bib2 Donald E. Knuth. {\tt mf.web}, version 2.71828182, 2021.
  \url{ctan.org\s tex-archive\s systems\s knuth\s dist\s mf}
\bib3 Donald E. Knuth. Adventure ({\tt advent.w}).
  \url{www-cs-faculty.stanford.edu\s \char`\~knuth\s programs\s
  advent.w.gz}
\bib4 Donald E. Knuth. The Stanford GraphBase.
  \url{www-cs-faculty.stanford.edu\s \char`\~knuth\s sgb.html}
\bib5 Donald E. Knuth. Literate programming.
  \url{www-cs-faculty.stanford.edu\s \char`\~knuth\s lp.html}
\bib6 Soojin Nam. \GWEB: Literate programming for Go.
  \url{github.com\s sjnam\s gweb}
\bib7 Soojin Nam. \GWEB\ examples.
  \url{github.com\s sjnam\s gweb-examples}
\bib8 Soojin Nam. metafont-the-program.
  \url{github.com\s sjnam\s metafont-the-program}
\bib9 Soojin Nam. tex-the-program.
  \url{github.com\s sjnam\s tex-the-program}
\bib{10} Dohyun Kim. \url{github.com\s dohyunkim}
\par}

\makesignature
\endarticle
